\section{Data}\label{sec:data}
% Budget: 2 pages. Generated from docs/data_notes.md and docs/decisions.md (section 1).

\paragraph{Universe and sample.} We use US-listed, USD-denominated country exchange-traded funds (ETFs) as tradable proxies for 17 emerging equity markets: Argentina (ARGT), Chile (ECH), Brazil (EWZ), Mexico (EWW) and Peru (EPU) in Latin America; South Korea (EWY), Taiwan (EWT), India (EPI), China (FXI), Thailand (THD), Indonesia (EIDO), Malaysia (EWM), Vietnam (VNM) and the Philippines (EPHE) in Asia; South Africa (EZA), Turkey (TUR) and Poland (EPOL) in EMEA. The S\&P~500 ETF (SPY) is included as a global reference node. Funds were selected if their price history starts by 3 March 2011 (7-day tolerance), at most 5\% of weeks are missing, and average daily traded volume is at least USD~3 million. Colombian funds (ICOL, GXG), INDA and MCHI did not pass the criteria; India and China are represented by EPI and FXI. The common sample runs from 11 March 2011 to 23 October 2026, i.e.\ 816 weeks, with no missing values.

\paragraph{Returns and frequency.} Prices are dividend- and split-adjusted, so returns are total returns. We use weekly Friday-to-Friday log returns. The weekly frequency follows the course indication that weekly returns retain some forecastability that is absent at daily frequency; it also mitigates non-synchronous trading between Asian and EMEA local markets and their US-listed ETFs. Unlike \citet{diebold2014network}, who study (realized) volatilities, we study returns: no daily data are available to build realized weekly volatility, and returns are the object of the forecastability argument. Returns are not devoid of cross-market predictive content: in a full-sample VAR(1) on the 18 series, 59 of the 324 slope coefficients (18.2\%) are significant at the 5\% level, against 5\% expected by chance. Our networks therefore describe the transmission of return shocks, not of risk.

\paragraph{Descriptive statistics.} Table~\ref{tab:descriptive} reports summary statistics. All series are strongly leptokurtic and mostly negatively skewed; an augmented Dickey--Fuller test rejects a unit root for every series. Five weekly log returns exceed 25\% in absolute value (ARGT in the week of the August 2019 primary elections and four observations in March 2020); they are genuine market moves and are kept. The March 2020 crash dominates the tails: of the ten largest weekly drops pooled across the 17 emerging markets, eight occur in the weeks ending 13 and 20 March 2020, and 14 of the 17 markets record their worst week in March 2020.

\begin{table}[ht]
\centering\scriptsize
\caption{Descriptive statistics of weekly log returns, 2011-03-11 to 2026-10-23 (816 weeks). Kurtosis is not in excess.}
\label{tab:descriptive}
\input{\tablepath/descriptive_stats.tex}
\end{table}

\paragraph{Limitations.} (i) ETF returns embed exchange-rate movements against the USD and are subject to non-synchronous trading. (ii) ARGT changed its underlying index in August 2014 (from a Solactive Argentina index to the MSCI All Argentina 25/50). (iii) FXI tracks only the 50 largest Chinese firms (FTSE China 50) and EPI is the earnings-weighted WisdomTree India fund, so neither is a broad market index. (iv) Colombia is absent. (v) Liquidity requirements and the common start date exclude younger funds. (vi) Without daily data, realized-volatility connectedness cannot be computed.
